\section{Herramientas y estrategia para las métricas de desempeño}
\label{sec:metricas}

\subsection{Escenarios y variables de la Tabla~1}

Cada escenario se ejecuta sobre la arquitectura base y sobre la escalable:

\begin{table}[H]
\centering
\caption{Escenarios de la Tabla~1.}
\begin{tabular}{lcccc}
\toprule
Escenario & Clientes & Solo consultan & Consultan y reservan & Reservan y modifican \\
\midrule
ESC-3 & 3 & 2 & 1 & 0 \\
ESC-5 & 5 & 3 & 0 & 2 \\
ESC-8 & 8 & 4 & 0 & 4 \\
\bottomrule
\end{tabular}
\end{table}

Se miden cuatro variables. Para efectos de la medición, la modificación cuenta
como operación de reserva y también se reporta por separado.

\begin{table}[H]
\centering
\caption{Variables y dónde se miden.}
\begin{tabularx}{\textwidth}{p{4.6cm}p{4.2cm}Y}
\toprule
Variable & Fórmula & Dónde se mide \\
\midrule
Tiempo promedio y desviación estándar de las \textbf{consultas} &
$\bar{t}=\frac{1}{n}\sum t_i$, \ $s=\sqrt{\frac{1}{n-1}\sum (t_i-\bar{t})^2}$ &
En cada cliente: inicio y fin de cada \codigo{consultarEventosMes} con el reloj
monótono (\codigo{std::chrono::steady\_clock}). \\
Tiempo promedio y desviación estándar de las \textbf{reservas} & Igual &
En cada cliente, para \codigo{reservar} y \codigo{modificarReserva}. \\
\textbf{TPM de consultas} & consultas exitosas $\div$ minutos de medición &
En el Gestor de Transacciones: contador por categoría y por minuto. \\
\textbf{TPM de reservas} & reservas confirmadas $\div$ minutos de medición &
En el Gestor y verificado en persistencia: transacciones de la bitácora con
instante dentro de la ventana de medición. \\
\bottomrule
\end{tabularx}
\end{table}

El tiempo de respuesta se mide en el cliente porque es lo que percibe el
estudiante: incluye red, gestor, procesador y persistencia. Las duraciones usan el
reloj monótono de una sola máquina, así que la deriva entre relojes no las afecta.
Las TPM se cuentan en el gestor de transacciones, como permite el enunciado, y se
cruzan con persistencia para confirmar que coinciden.

\subsection{Instrumentación en el código}

\begin{itemize}
    \item \textbf{Cliente de carga:} registra por operación el instante de inicio, la
          categoría (\codigo{CONSULTA} o \codigo{RESERVA}), la operación, la latencia
          en milisegundos, si fue exitosa y el código de error. Al terminar escribe un
          archivo CSV:
\begin{lstlisting}
escenario,arquitectura,cliente,tipo_cliente,inicio_ms,categoria,operacion,latencia_ms,exitosa,codigo_error
ESC-5,base,C1,consulta,1792706556363,CONSULTA,consultarEventosMes,4,1,
\end{lstlisting}
    \item \textbf{Gestor de Transacciones:} una clase \codigo{Metricas} acumula por
          categoría la cantidad, los errores, la suma y la suma de cuadrados de las
          latencias, y cada minuto agrega al archivo
          \codigo{metricas\_gestor.csv} una línea con transacciones, errores,
          promedio, desviación y TPM. Agregar un contador cuesta unos pocos
          nanosegundos, así que no altera lo que se mide.
    \item \textbf{Persistencia:} cada transacción guarda su instante en la bitácora,
          lo que permite contar las reservas confirmadas por minuto.
    \item \textbf{Recursos:} CPU y memoria por proceso con \codigo{pidstat} (paquete
          \codigo{sysstat}) cada segundo, para identificar el cuello de botella.
\end{itemize}

\subsection{Herramientas}

\begin{table}[H]
\centering
\caption{Herramientas y librerías de medición.}
\begin{tabularx}{\textwidth}{p{4.4cm}Y}
\toprule
Herramienta & Uso \\
\midrule
\codigo{std::chrono} (C++17) & Reloj monótono para las latencias y reloj de pared
para ubicar cada medición en la ventana. \\
Archivos CSV & Formato común de los clientes, el gestor y \codigo{pidstat}. \\
Python 3 + \textbf{pandas} y \textbf{NumPy} & Unir los CSV de todos los clientes,
descartar el calentamiento y calcular promedio, desviación estándar muestral, p95,
TPM y porcentaje de transacciones < 1\,s por escenario y arquitectura. \\
\textbf{matplotlib} & Gráficos para el informe. \\
\codigo{pidstat} / \codigo{htop} & Uso de CPU y memoria por proceso. \\
Prometheus + Grafana (segunda entrega) & Paneles en tiempo real con la librería
\textbf{prometheus-cpp} en los procesos, y \codigo{node\_exporter} para las
máquinas. Útiles para las pruebas de estrés y de fallos. \\
\bottomrule
\end{tabularx}
\end{table}

\subsection{Procedimiento de cada corrida}

\begin{enumerate}
    \item Levantar la arquitectura (base: 1 procesador RMC; escalable: 3) con datos
          iniciales limpios.
    \item Lanzar a la vez los procesos cliente del escenario.
    \item Descartar los primeros 2 minutos (calentamiento) y medir durante 10 minutos.
    \item Repetir cada combinación 3 veces y reportar el promedio de las repeticiones.
    \item Verificar la consistencia (C-03): ninguna fecha con sobrecupo.
\end{enumerate}

Son 3 escenarios $\times$ 2 arquitecturas $\times$ 3 repeticiones = 18 corridas.

\subsection{Presentación de los resultados}

Los resultados se presentarán en una tabla con una fila por escenario y
arquitectura (promedio y desviación de consultas y de reservas, TPM de consultas y
de reservas, errores) y en cuatro gráficos con el número de clientes (3, 5 y 8) en
el eje horizontal y una línea por arquitectura:

\begin{enumerate}
    \item tiempo promedio de las consultas, con barras de $\pm$ una desviación estándar;
    \item tiempo promedio de las reservas, con barras de $\pm$ una desviación estándar;
    \item TPM de consultas;
    \item TPM de reservas.
\end{enumerate}

\begin{table}[H]
\centering
\caption{Plantilla de la tabla de resultados (se completará en la entrega final).}
\small
\begin{tabular}{llcccccc}
\toprule
& & \multicolumn{2}{c}{Consultas (ms)} & \multicolumn{2}{c}{Reservas (ms)} & \multicolumn{2}{c}{TPM} \\
\cmidrule(lr){3-4}\cmidrule(lr){5-6}\cmidrule(lr){7-8}
Escenario & Arquitectura & Prom. & Desv. & Prom. & Desv. & Consultas & Reservas \\
\midrule
ESC-3 & Base & & & & & & \\
ESC-3 & Escalable & & & & & & \\
ESC-5 & Base & & & & & & \\
ESC-5 & Escalable & & & & & & \\
ESC-8 & Base & & & & & & \\
ESC-8 & Escalable & & & & & & \\
\bottomrule
\end{tabular}
\end{table}
